% ---
% id: meat-core/frank-wolfe/rem:iterate-relaxed-oracle-exact
% kind: rem
% title: "Iterate in the relaxed set, minima over the exact set"
% summary: >-
%   A run establishes its iterate only in the relaxed feasible set, while the
%   linear minima its bracket bounds are over the exact one, and part (iii) of
%   [[thm:firstorder]] asks each half of the object that supplies it. Its bound
%   compares h_lb with the exact h_low directly, through no relaxed quantity.
% prominence: standard
% status: accepted
% uses: [thm:firstorder, def:first-order-setting, prop:relaxed-set-properties]
% mentions: [rem:feas-tol, lem:feasible-side-dual, rem:acc-side, cor:fo-keylength, meat-core/thm:choi, meat-core/prop:value-range, foundations/def:exact-island, req:iterate-checks]
% assumes: []
% refs: []
% provenance: newly-derived
% lean:
% ---

\begin{remark}\label{rem:iterate-relaxed-oracle-exact}
Part~(iii) of Theorem~\ref{thm:firstorder} substitutes the
\tref{relaxed feasible set} $\Feas^{J}(\varepsilon)$ in different places
in its two halves, and the asymmetry is deliberate. A run of this module
is always in that part: it verifies the $J$ factor of its iterate's
membership at positive declared windows
(Requirement~\ref{req:iterate-checks}) and holds the $\acSet$ factor
exactly, by forming $\mathbf q_t$ on the \tref{exact island}
(\xref{foundations/def:exact-island}; Remark~\ref{rem:acc-side}), so what
it establishes is $x_t\in\acSet\times\Feas^{J}(\varepsilon)$ and not exact
membership in the product; while the minima of \eqref{eq:lmo}, and the
lower bounds that part~(iv) of the theorem puts in their place, are taken
over the declared conic program \eqref{eq:choi-feas}, which is the
\emph{exact} $\Feas^{J}$. \emph{Iterate in the relaxed set, minima over
the exact set}: (iiib) asks each half of the object that supplies it. Its
hypotheses (H-POS), (H-SUPP) and (H-NOPERT) are conditions on the iterate
alone and name no set. Substituting $\Feas^{J}(\varepsilon)$ into
\eqref{eq:lmo} as well would ask for a bound against a \emph{larger} set,
hence below a \emph{smaller} minimum: since
$\Feas^{J}\subseteq\Feas^{J}(\varepsilon)$, one has
$\min_{\Feas^{J}(\varepsilon)}\le\min_{\Feas^{J}}$, so the inequality over
the exact set is the weaker of the two. It is also the one
Lemma~\ref{lem:feasible-side-dual} supplies: the trace identity that lemma
rests on is the exact coupling condition, which a point of the relaxed set
meets only within its window.

Part~(iiib) concludes against the exact $\hlow$ directly, with no
intermediate quantity: $\hlow[\Feas^{J}(\varepsilon)]$ appears in (iiia)
and nowhere in (iiib), and no non-negativity of it is claimed; it may be
negative, by at most $\frac{\alpha}{\alpha-1}\log_2(1+d_{A_0}\varepsilon)$
(Proposition~\ref{prop:relaxed-set-properties}\,(c$''$)). The clamp at
zero reads the range fact $\hlow\ge0$ of
Proposition~\ref{prop:value-range}\,(i) on the exact set, which is the only
reading it needs; on the relaxed set that fact is not available, and
nothing uses it there. Corollary~\ref{cor:fo-keylength} consumes the same
inequality on this reading as on that of~(ii), (ii) being the case
$\varepsilon=0$ of (iiib). Which $\varepsilon$ a run uses, the larger of
its two declared windows on computed residuals and never the declared
tolerance, is fixed by Requirement~\ref{req:iterate-checks}, and part~(iii)
is insensitive to it, being proved for every $\varepsilon\ge0$.

The relaxed quantity in (iiia) is the Choi-form one and not the state-form
$\hlow[\Feas(\varepsilon)]$ of
Proposition~\ref{prop:relaxed-set-properties}\,(d): at $\varepsilon>0$ the
two relaxed sets are not in bijection under $\chi_{\Qsys}$, for the two
reasons Remark~\ref{rem:feas-tol} gives, and no relation between the two
numbers is claimed or needed. Part~(iiib) uses neither relaxed value: it
compares $\hlb$ against the exact $\hlow$ directly, so the chain
\eqref{eq:feas-tol-mono-J} is consumed by the estimate half of part~(iii)
and not by the bound half.
\end{remark}
